\subsection*{Exercises \thesection}

% counterexample: convergence in distribution but not in probability

\begin{exercise}[Convergence in distribution but not probability \Coffeecup]
    \label{ex: convergence in distribution but not in probability}
    Let \(X\) be a random variable with \(\prob{X=-1} = \frac12\) and \(\prob{X=1}=\frac12\).
    \[
        X_n \coloneq (-1)^n X
    \]
    Prove that
    \begin{enumerate}[label=(\alph*), noitemsep]
        \item \(X_n \xrightarrow{d} X\)
        \begin{solution}
            Observe that the distribution of \(X\) is symmetric, that is
            \[
                \prob{-X \le x} = \prob{X \le x} \qquad \forall x\in \real.
            \]
            Consequently, for any \(x \in \real\) we have \(\prob{X_n \le x} = \prob{X \le x}\).
            In particular this implies \(\prob{X_n \le x} \to \prob{X \le x}\) for all \(x\).
        \end{solution}

        \item \(X_n\) does not converge in probability to \(X\).
        \begin{solution}
            We have for \(\epsilon = 1\) and every odd \(n\)
            \[
                \prob{\abs{X - X_n} \ge \epsilon}
                = \prob{2\abs{X} \ge 1} = 1
            \]
            The sequence \(\prob{|X_n - X| \ge \epsilon}\) thus does not converge to \(0\).
        \end{solution}
    \end{enumerate}
\end{exercise}



\begin{exercise}[Convergence in probability but not \(L^p\) \Coffeecup]
    \label{ex: convergence in probability but not in Lp}
    Let \(U\sim \uniform(0,1)\) be a standard uniform random variable and define
    \[
        X_n \coloneq n \ind{[0, 1/n)}(U)
    \]
    Show that
    \begin{enumerate}[label=(\alph*), noitemsep]
        \item the sequence \((X_n)_{n\in \nat}\) converges to \(0\) in probability,
        \begin{solution}
            We have for any \(\epsilon > 0\)
            \[
                \prob{|X_n - 0| \ge \epsilon}
                \le \prob(X_n \neq 0)
                = \prob[\big]{U\le \tfrac1n}
                = \tfrac1n
                \to 0.
                \qedhere
            \]
        \end{solution}

        \item the sequence \((X_n)_{n\in \nat}\) does not converge to \(0\) in \(L^1\),
        \begin{solution}
            We have
            \[
                \E{|X_n - 0|}
                = \E{X_n}
                = n \cdot \prob[\big]{U \in [0, 1/n)}
                = n \cdot \frac1n
                = 1 \not\to 0.
                \qedhere
            \]            
        \end{solution}

        \item For every \(p \ge 1\) construct a sequence \((Y_n)_{n\in \nat}\) such that
        \(Y_n\xrightarrow{L^q} 0\) for all \(q < p\) but not for \(q \ge p\).
        \begin{solution}
            Define \(Y_n \coloneq X_n^{1/p}\) with \(X_n\) as above. Then
            \[
                \E{|Y_n - 0|^q} = \E{X_n^{q/p}} = n^{q/p} \cdot \prob[\big]{U \in [0, 1/n)} = n^{q/p - 1}.
            \]
            This converges to \(0\) for all \(q < p\) but not for \(q \ge p\).
        \end{solution}
        \end{enumerate}
\end{exercise}


\begin{exercise}[Convergence in \(L^p\) but not almost surely \Coffeecup\Coffeecup] 
    \label{ex: convergence in Lp but not as}
    Let \(U\sim \uniform(0,1)\) be a standard uniform random variable.
    We will scan the interval for \(U\) with increasingly fine measurements. Specifically,
    let
    \begin{alignat*}{7}
        X_1 &\coloneq \ind{[0,1)}(U)
        \\
        X_2 &\coloneq \ind{[0,\frac12)}(U)
        \quad & X_3 &\coloneq \ind{[\frac12, 1)}(U)
        \\
        X_4 &\coloneq \ind{[0, \frac14)}(U)
        \quad & X_5 &\coloneq \ind{[\frac14, \frac12)}(U)
        \quad & X_6 &\coloneq \ind{[\frac12, \frac34)}(U)
        \quad & X_7 &\coloneq \ind{[\frac34, 1)}(U)
    \end{alignat*}
    That is
    \[
        X_n \coloneq \ind{\left[
            \frac{k}{2^{m}}, \frac{k+1}{2^{m}}
        \right)}(U),
        \quad \text{with}\quad 
        m \coloneq \lfloor\log_2(n)\rfloor, \quad k \coloneq n - 2^m. 
    \]
    Show that
    \begin{enumerate}[label=(\alph*), noitemsep]
        \item\label{it: (conv in Lp but not as) in Lp}
        \(X_n \xrightarrow{L^p} 0\),

        \item\label{it: (conv in Lp but not as) in prob}
        \(X_n \xrightarrow{p} 0\),

        \item\label{it: (conv in Lp but not as) not as}
        \(X_n\) does not converge almost surely to \(0\).
    \end{enumerate}
\end{exercise}
\begin{solution}
\begin{enumerate}[wide=0pt]
    \item[\ref{it: (conv in Lp but not as) in Lp}]
    We have for any \(p \ge 1\)
    \[
        \E{|X_n|^p} = \int_0^1 \ind{
            \left[\frac{k}{2^{m}}, \frac{k+1}{2^{m}} \right)
        }(u) \, du
        = \frac{1}{2^m} = 2^{-\lfloor \log_2(n) \rfloor} \to 0.
    \]
    Consequently, \(X_n \xrightarrow{L^p} 0\).

    \item[\ref{it: (conv in Lp but not as) in prob}]
    \(X_n \xrightarrow{L^p} 0\) implies \(X_n \xrightarrow{p} 0\) by Theorem
    \ref{thm: convergence implications} \ref{it: L1 to p}.

    Alternatively:
    \[
        \prob[\bigl]{|X_n - 0| \ge \epsilon}
        = \prob{X_n = 1}
        = 2^{-\lfloor \log_2(n) \rfloor} \to 0.
    \]
    for any \(\epsilon \in (0,1]\) and \(\prob[\bigl]{|X_n - 0| \ge
    \epsilon}=0\) for \(\epsilon > 1\).

    \item[\ref{it: (conv in Lp but not as) not as}]
    Choose any \(\omega \in \set{U\in [0,1)}\). We will show that there exists an
    \(\epsilon > 0\), specifically any \(\epsilon < 1\) such that
    \(|X_n(\omega) - 0| \ge \epsilon\) for infinitely many \(n\).
    Since we show this for any \(\omega\in \set{U\in [0,1)}\), this would imply
    \[
        \prob{X_n \not\to 0} \ge \prob{U \in [0,1)} = 1.
    \]
    For any \(N\in \nat\) the task is consequently to find \(n\ge N\)
    such that \(X_n(\omega) = 1\). Let \(m \coloneq \lfloor \log_2(N) \rfloor + 1\).
    Then there exists a \(k \in \{0, \ldots, 2^m-1\}\) such that
    \[
        \frac{k}{2^m} \le U(\omega) < \frac{k+1}{2^m}.
    \]
    Consequently, for \(n = 2^m + k\) we have \(n \ge N\) and
    \[
        X_n(\omega) = \ind{\left[
            \frac{k}{2^{m}}, \frac{k+1}{2^{m}}
        \right)}(U(\omega)) = 1.
        \qedhere
    \]
\end{enumerate}
\end{solution}


\begin{exercise}[subtitle={Stochastic convergence sometimes implies \(L^p\) \Coffeecup}]
    \label{ex: p to Lp when bounded}
    Assume that \(X_n\) and \(X\) are bounded random variables. That is 
    \(\sup_n\|X_n\| \le c\) and \(\|X\|\le c\) with probability one for some \(c
    \in \real\). Show that then \(X_n \overset{p}\to X\) implies \(X_n
    \overset{L^p}\to X\).
    \begin{remark}
        This is a weak form of the dominated convergence theorem, where the constant
        is the dominating function.
    \end{remark}
\end{exercise}
\begin{solution}
    We have for every \(\epsilon > 0\)
    \[\begin{aligned}
        \E[\big]{\norm{X_n - X}^p}
        &= \E[\big]{\norm{X_n - X}^p \ind{\norm{X_n - X} < \epsilon}}
        + \E[\big]{\norm{X_n - X}^p \ind{\norm{X_n - X} \ge \epsilon}}
        \\
        &\le \epsilon^p + (2c)^p \prob[\big]{\norm{X_n - X} \ge \epsilon}.
    \end{aligned}
    \]
    Due to \(X_n \overset{p}\to X\) the second term converges to \(0\)
    and we thus have
    \[
        \limsup_{n\to \infty} \E[\big]{\norm{X_n - X}^p} \le \epsilon^p.
    \]
    Since this is true for any \(\epsilon > 0\) we have the claim.
\end{solution}


\begin{exercise}[Stochastic triangle inequality \Coffeecup]
    \label{ex: stochastic triangle}
    Let \(X,Y,Z\) be random variables. Show for all \(\epsilon > 0\)
    \[
        \prob*{\norm{X - Z} \ge \epsilon} \le \prob*{\norm{X - Y} \ge \tfrac\epsilon2} + \prob*{\norm{Y - Z} \ge \tfrac\epsilon2}.
    \] 
    \textbf{Hint:} Use set inclusions.
    \begin{solution}
        The claim follows directly from
        \[
            \set{\norm{X - Z} \ge \epsilon}
            \subseteq \set{\norm{X - Y} \ge \tfrac\epsilon2} \cup \set{\norm{Y - Z} \ge \tfrac\epsilon2},
        \]
        which is proven using that \(A \subseteq B\) is equivalent
        to \(B^\complement \subseteq A^\complement\) and
        \begin{align*}
            \Bigl(\set{\norm{X - Y} \ge \tfrac\epsilon2} \cup \set{\norm{Y - Z} \ge \tfrac\epsilon2}\Bigr)^\complement
            &=\set{\norm{X - Y} < \tfrac\epsilon2} \cap \set{\norm{Y - Z} < \tfrac\epsilon2}
            \\
            &\subseteq \set{\norm{X - Z} < \epsilon}
            \\
            &= \set{\norm{X - Z} \ge \epsilon}^\complement.
        \end{align*}
        The standard triangle inequality is used for the inclusion step.
    \end{solution}
\end{exercise}


\begin{exercise}[Stochastic limits are unique \Coffeecup\Coffeecup]
    \label{ex: unique stochastic limits}
    Show that
    \begin{enumerate}[label=(\alph*),noitemsep]
        \item if \(X_n \overset{p}\to X\) and \(X_n \overset{p}\to Y\), then
        \(\prob{X = Y}=1\).

        \textbf{Hint:} How is uniqueness of limits proven in analysis?
        \begin{solution}
            We have for any \(\epsilon > 0\) and any \(n\in \nat\) by the
            stochastic triangle inequality (Exercise \ref{ex: stochastic
            triangle})
            \[\begin{aligned}
                \prob{\norm{X - Y} \ge \epsilon}
                \le \underbrace{\prob{\norm{X - X_n} \ge \tfrac\epsilon2}}_{\to 0} + \underbrace{\prob{\norm{X_n - Y} \ge \tfrac\epsilon2}}_{\to 0},
            \end{aligned}
            \]
            Since this inequality holds for any \(n\in \nat\) and \(X_n\)
            converges to both \(X\) and \(Y\) in probability, this term must
            therefore be equal to zero. That is \(\prob{\norm{X - Y} \ge
            \epsilon} = 0\) for any \(\epsilon > 0\). Since this is true for any
            \(\epsilon > 0\) we have by continuity of the probability measure
            \[
                \prob{X \neq Y} = \prob[\Big]{\bigcup_{n\in \nat}\norm{X - Y} > \frac1n} = \lim_{n\to\infty} \prob{\norm{X - Y} \ge \tfrac1n} = 0.
            \]
        \end{solution}
    \end{enumerate}
    Deduce
    \begin{enumerate}[label=(\alph*), resume,noitemsep]
        \item if \(X_n \overset{\as}\to X\) and \(X_n \overset{\as}\to Y\), then \(\prob{X = Y}=1,\)
        \item if \(X_n \overset{L^p}\to X\) and \(X_n \overset{L^p}\to Y\), then \(\prob{X = Y}=1.\)
        \begin{solution}
            If \(X_n \overset{\as}\to X\) and \(X_n \overset{\as}\to Y\), then
            \(X_n \overset{p}\to X\) and \(X_n \overset{p}\to Y\) by Theorem
            \ref{thm: convergence implications} \ref{it: as to p}. The claim
            then follows from the previous part. Similarly for \(L^p\) convergence.
        \end{solution}
    \end{enumerate}
    Finally, 
    \begin{enumerate}[label=(\alph*), resume, noitemsep]
        \item\label{it: non-unique limits convergence in dist} show that the statement is not true for convergence in
        distribution.

        \textbf{Hint:} Clearly \(X_n\) must converge in distribution but not
        in probability as the limits would otherwise be almost surely equal by the
        previous parts.

        \begin{solution}
            We may choose \(X\) as in Exercise \ref{ex: convergence in distribution but not in probability}.
            Since \(X_n \coloneq X\) converges both to \(X\) and \(-X\) in
            distribution the limit is clearly not unique as
            \[
                \prob{X = -X} = \prob{X=0}= 0.
                \qedhere
            \]
        \end{solution}
    \end{enumerate}
\end{exercise}

\begin{exercise}[subtitle={Stochastic Simulation and Quantiles \Coffeecup\Coffeecup}]
    \label{ex: stochastic simulation}
    The goal is to simulate a random variable \(X\) with distribution 
    function \(F(x) = \prob{X \le x}\) using a standard uniform random variable
    \(U\sim \uniform(0,1)\). We define the generalized inverse of \(F\) as
    \[
        F^{\leftarrow}(p) \coloneq \inf\{x \in \real : F(x) \ge p\}, \quad p \in (0,1).
    \]
    Show that
    \begin{enumerate}[label=(\alph*), noitemsep]
        \item\label{it: quantile if invertible}
        if \(F\) is invertible, then \(F^{\leftarrow}(p) = F^{-1}(p)\) for all
        \(p \in (0,1)\),

        \item\label{it: monotonicity}
        \(F^{\leftarrow}\) is non-decreasing,

        \item\label{it: F-(F(x)) le x}
        \(F^{\leftarrow}(F(x)) \le x\) for all \(x \in \real\),

        \item\label{it: F(F-(p)) > p}
        \(F(F^{\leftarrow}(p)) \ge p\) for all \(p \in (0,1)\),

        \item\label{it: F-(p)< x iff p<F(x)} \(F^{\leftarrow}(p) \le x\) if and only if \(p \le F(x)\),

        \item\label{it: inverse transform sampling}
        \(X\coloneq F^{\leftarrow}(U) \sim F\) for a standard uniform \(U\sim \uniform(0,1)\), i.e.\ 
        \[
            F(x) = \prob{X \le x}
            \qquad \forall x \in \real.
        \]

        \item\label{it: simulate exponential rvs} Propose a way to simulate an exponential distribution \(X \sim
        \exponential(\lambda)\) using a standard uniform random variable \(U\sim
        \uniform(0,1)\).
    \end{enumerate}
    The function \(F^{\leftarrow}\) is also called a \emph{quantile function} of
    the distribution function \(F\). The reason for this is
    the following: Show that
    \begin{enumerate}[resume, label=(\alph*), noitemsep]
        \item\label{it: quantiles}
         \(Q(p) \coloneq (-\infty, F^{\leftarrow}(p)]\) is the smallest
        interval of the form \((-\infty, a]\) such that \(\prob{X \in (-\infty,
        a]} \ge p\) for \(X\sim F\). We say that \(Q(p)\) is a \(p\)-quantile.
    \end{enumerate}
\end{exercise}
\begin{solution}
\begin{enumerate}[wide=0pt]
    \item[\ref{it: quantile if invertible}]
    Since \(F\) is monotonically increasing as a distribution
    function and additionally invertible, it is strict monotonically
    increasing. That is
    \[
        F(x) < F(y) \quad \text{for all } x < y.
    \]
    So clearly for all \(x< F^{-1}(p)\) we have \(F(x) < p\) and
    therefore they are not contained in the set \(\set{x\in \real : F(x)
    \ge p}\). This implies \(F^{-1}(p)\) is the smallest element in this set
    and consequently \(F^{\leftarrow}(p) = F^{-1}(p)\).

    \item[\ref{it: monotonicity}]
    Let \(p < q\). Then for any \(x \in \real\) such that \(F(x) \ge q\)
    we also have \(F(x) \ge p\). Consequently,
    \[
        \set{x \in \real : F(x) \ge q} \subseteq \set{x \in \real : F(x) \ge p},
    \]
    This implies the infimum of the former is greater than the infimum of the latter
    and thus
    \[
        F^{\leftarrow}(p) = \inf\{x \in \real : F(x) \ge p\} \le \inf\{x \in \real : F(x) \ge q\} = F^{\leftarrow}(q).
    \]

    \item[\ref{it: F-(F(x)) le x}]
    Since \(x\in \set{y\in \real : F(y) \ge F(x)}\), we have
    \[
        x \ge \inf\set{y\in \real : F(y) \ge F(x)} = F^{\leftarrow}(F(x)).
    \]

    \item[\ref{it: F(F-(p)) > p}]
    Let \(x_n \in \set{x\in \real : F(x) \ge p}\) be a sequence with \(x_n \downarrow F^{\leftarrow}(p)\). Then
    by right-continuity of the distribution function \(F\)
    \[
        p \le \lim_{n\to \infty} F(x_n)
        = F\Bigl(\lim_{n\to \infty} x_n\Bigr)
        = F(F^{\leftarrow}(p)).
    \]
    
    \item[\ref{it: F-(p)< x iff p<F(x)}]
    Assume \(F^{\leftarrow}(p) \le x\). Then by \ref{it: F(F-(p)) > p} 
    and the monotonicity of \(F\) we have
    \[
        p \overset{\ref{it: F(F-(p)) > p}}\le F(F^{\leftarrow}(p)) \le F(x).
    \]
    For the opposite direction assume \(p\le F(x)\), then by the
    monotonicity of \(F^{\leftarrow}\) \ref{it: monotonicity}
    and \ref{it: F-(F(x)) le x} we have
    \[
        F^{\leftarrow}(p) \overset{\ref{it: monotonicity}}\le F^{\leftarrow}(F(x)) \overset{\ref{it: F-(F(x)) le x}}\le x.
    \]

    \item[\ref{it: inverse transform sampling}]
    We simply use \ref{it: F-(p)< x iff p<F(x)} and the definition of \(U\) to get
    \[
        \prob{X \le x}
        = \prob{F^{\leftarrow}(U) \le x}
        \overset{\ref{it: F-(p)< x iff p<F(x)}}\le \prob{U \le F(x)}
        = \int_0^{F(x)} 1 \, du
        = F(x).
    \]

    \item[\ref{it: simulate exponential rvs}]
    The CDF of an exponential distribution is \(F(x) = 1 - e^{-\lambda x}\) for \(x \ge 0\).
    Since this function is invertible on \([0, \infty]\) we use
    \ref{it: quantile if invertible}. To get
    \[
        F^{\leftarrow}(p) = -\tfrac{1}{\lambda} \log(1-p).
    \]
    Using \(U \sim \uniform(0,1)\) we thus set \(X = F^{\leftarrow}(U)\),
    and finish by \ref{it: inverse transform sampling}.

    \item[\ref{it: quantiles}]
    For any \(a < F^{\leftarrow}(p)\) we have by definition of \(F^{\leftarrow}\) that
    \(F(a) < p\) and thus
    \[
        \prob{X \in (-\infty, a]} = F(a) < p.
    \]
    However for \(a=F^{\leftarrow}(p)\) we have \(F(a)\ge p\) by
    \ref{it: F(F-(p)) > p}.  Consequently, \(F^{\leftarrow}(p)\) is the
    smallest \(a\) such that \(\prob{X \in (-\infty, a]} \ge p\).
    \qedhere
\end{enumerate}
\end{solution}

\begin{exercise}[subtitle={Monotone functions have countable discontinuities \Coffeecup\Coffeecup\Coffeecup}]
    \label{ex: monotone functions have countably many discontinuities}
    Let \(F\colon \real \to \real\) be a non-decreasing function.
    \begin{enumerate}[label=(\alph*), noitemsep]
        \item\label{it: (monotone functions) left and right limits}
        Let \(x\in \real\). Prove that the limit
        \[
            F_-(x) \coloneq \lim_{n\to \infty} F(x_n)
        \]
        exists for any sequence \((x_n)_{n\in \nat} \subseteq \real\) with \(x_n < x\)
        and \(x_n \to x\). Prove that this limit is the same for any such
        sequence such that \(F_-(x)\) is well-defined.
        Also prove that \(F_+(x) \coloneq \lim_{y\downarrow x} F(y)\) is
        well-defined.

        \textbf{Hint:} Consider monotonically increasing sequences \(x_n\) first.
        Then prove that the liminf and limsup coincide for general sequences
        using monotone subsequences.

        \item\label{it: F continuous if and only if all coincide}
        Prove that \(F\) is continuous at \(x\) if and only
        if \(F(x) = F_-(x) = F_+(x)\). That is
        \(F\) only has jump discontinuities.

        \item\label{it: (monotone functions) countable jumps, compact interval}
        Prove that \(F\) restricted to a compact interval \([a,b]\)
        has only countably many jumps (places where \(F_-(x) < F_+(x)\)).

        \textbf{Hint:} How many jumps are there of size greater than \(1/n\)?

        \item\label{it: (monotone functions) countable jumps, real line}
        Prove that \(F\) has at most countably many discontinuities.
    \end{enumerate}
\end{exercise}

\begin{solution}
\begin{enumerate}[wide=0pt]
    \item[\ref{it: (monotone functions) left and right limits}]
    Let \(x_n\uparrow x\) be a monotonically increasing sequence. Then
    \(F(x_n)\) is also monotonically increasing and thus assumes a limit.
    We will prove \(\lim_n F(x_n) \le \lim_n F(y_n)\)
    for another monotonically increasing sequence \(y_n\uparrow x\).
    Since \(x_n\) and \(y_n\) are exchangeable we then also have equality.
    Now for any \(n\in \nat\) there exists an \(m_n\in \nat\) such that
    \(x_n \le y_{m_n} \le x\) as \(x_n < x\) and \(y_n\to x\). Consequently
    since \(F\) is non-decreasing we have
    \[
        \lim_{n\to\infty} F(x_n) \le \lim_{n\to\infty}F(y_{m_n}) \overset{\text{subsequence}}= \lim_{n\to\infty}F(y_n).
    \]
    Thus we have proven that the limit of monotonically increasing
    sequences is unique and denoted by \(F_-(x)\). Let \(x_n \to x\) with \(x_n < x\) be an
    arbitrary sequence. Let \(x_{n_k}\) be a subsequence with
    \[
        F(x_{n_k}) \to \liminf_{n\to\infty} F(x_n)
    \]
    Since \(x_{n_k} \to x\) and \(x_{n_k} < x\) we can find a monotonically increasing subsequence.
    This subsequence converges to \(F_-(x)\) and since it must have the
    same limit as the original sequence \(x_{n_k}\) we have
    \[
        \liminf_{n\to\infty} F(x_n) = F_-(x).       
    \]
    The same argument works for the limsup. Thus the limit exists and is
    given by \(F_-(x)\). The exact same argument also works for
    \(F_+(x)\) with decreasing sequences.

    \item[\ref{it: F continuous if and only if all coincide}]
    Clearly we have \(F_-(x) = F_+(x)\) if \(F\) is continuous.  So we
    only need to prove \(F\) continuous in \(x\) when this equality
    holds true.  Note that
    \(F_-(x) \le F(x) \le F_+(x)\) since \(F\) is non-decreasing.
    So if we have the equality then by a sandwich argument they are
    both equal to \(F(x)\). Let us assume equality. To prove
    sequential continuity in \(x\) we choose a sequence \(x_n\to x\).
    This sequence can be split into the parts that are smaller than \(x\)
    converging to \(F_-(x)\), the parts that are equal to \(x\),
    converging to \(F(x)\) and the parts that are larger than \(x\)
    converging to \(F_+(x)\). Since all these limits coincide the entire
    sequence converges to \(F(x)\). Consequently \(F\) is continuous in
    \(x\).

    \item[\ref{it: (monotone functions) countable jumps, compact interval}]
    Since \(F(b) - F(a)\) is finite and \(F\) monotonically increasing,
    there are at most \(\lceil n(F(b)-F(a))\rceil\) jumps of size
    \(\frac1n\). Thus for any \(n\in \nat\) the number of jumps of greater
    size is finite. This implies that the discontinuity set
    \[
        D_F = \set{x: F_-(x) \neq F_+(x)} = \bigcup_{n\in \nat} \set[\big]{x: F_+(x) - F_-(x) > \tfrac1n}
    \]
    is countable.

    \item[\ref{it: (monotone functions) countable jumps, real line}]
    Since \(F\) only has countably many discontinuities in each
    interval \([n, n+1]\) for \(n\in \integer\) and the countable union
    of countable sets is countable, we have that the set of all
    discontinuities is at most countable.
\end{enumerate}
\end{solution}


\begin{exercise}[Discontinuity points of distribution functions \Coffeecup\Coffeecup]
    Let \(X\) be a random variable and \(F(x) = \prob{X \le x}\) its
    distribution function. Show that \(F\) is continuous at \(x\) if and only if
    \(\prob{X = x} = 0\).
    \textbf{Hint:} Use Exercise \ref{ex: monotone functions have countably many discontinuities}.
\end{exercise}
\begin{solution}
    Since the sets \(\set{X \le x - \frac1n}\) are increasing
    in size, we have by continuity of the probability measure
    \[
        \prob{X < x}
        = \prob[\Big]{\bigcup_{n\in \nat} \set{X \le x - \tfrac1n}}
        = \lim_{n\to\infty} \prob{X \le x - \tfrac1n}
        = F_-(x).
    \]
    Consequently,
    \[
        \prob{X = x} = \prob{X \le x} - \prob{X < x} = F(x) - F_-(x).
    \]
    By Exercise \ref{ex: monotone functions have countably many
    discontinuities} \ref{it: F continuous if and only if all coincide} we
    have that \(F\) is continuous at \(x\) if and only if \(F(x) = F_-(x)\).
    Consequently continuity of \(F\) is equivalent to \(\prob{X = x} = 0\).
\end{solution}

\begin{exercise}[Skorokhod's representation theorem \Coffeecup\Coffeecup]
    \label{ex: skorokhod representation}
    Let \((X_n)_{n \in \nat}\) be a sequence of random variables in \(\real\)
    such that \(X_n \xrightarrow{d} X\). Show that there exist random variables
    \((Y_n)_{n\in \nat}\) and \(Y\), where the \(Y_n\) and \(Y\) are all defined on the
    same probability space such that
    \begin{enumerate}[label=(\roman*), noitemsep]
        \item \(Y_n \overset{d}= X_n\) for all \(n\in \nat\) and \(Y \overset{d}= X\),
        \item \(Y_n \overset{\as}\to Y\).
    \end{enumerate}
    \textbf{Hint:} Use \(F_n(x) = \prob{X_n \le x}\), Exercise \ref{ex:
    stochastic simulation} and that a monotonous function has at most countably
    many discontinuities (Exercise \ref{ex: monotone functions have countably
    many discontinuities}).
    \begin{remark}
        We assume \(X_n\) in \(\real\) to simplify the proof. Skorokhod's representation 
        theorem holds in much greater generality.
    \end{remark}
\end{exercise}
\begin{solution}
    Choose \(U\sim \uniform(0,1)\) and define
    \[
        Y_n \coloneq F_n^{\leftarrow}(U)
        \quad \text{and}\quad
        Y = F^{\leftarrow}(U)
    \] where \(F_n\) are the distribution functions of \(X_n\) and \(F\) the
    one of \(X\). Then by Exercise \ref{ex: stochastic simulation} we have
    \(Y_n \overset{d}= X_n\) for all \(n\in \nat\) and
    \(Y \overset{d}= X\).  It remains to show that \(Y_n \overset{\as}\to
    Y\). Since \(X_n \overset{d}\to X\) in distribution, we have that
    \(F_n(x) \to F(x)\) for all \(x\in \real\) at which \(F\) is continuous.
    Since the set \(D_F\) of discontinuity points of \(F\) is countable, we have
    \[
        \prob{U \in D_F} = \sum_{x\in D_F}\prob{U = x} = 0.
    \]
    And for all \(\omega \in \set{U \notin D_F}\) we have
    \[
        Y_n(\omega) = F_n(U(\omega)) \to F(U(\omega))=Y(\omega)
    \]
    Consequently, \(\prob{ Y_n \to Y } = \prob{U \notin D_F} = 1\) and thus \(Y_n \overset{\as}\to Y\).
\end{solution}


\begin{exercise}[Lipschitz continuous bump functions \Coffeecup\Coffeecup]
    \label{ex: lipschitz bump functions}
    Let \((\domain, \metric)\) be a metric space and \(A \subseteq \domain\). 
    Show that
    \begin{enumerate}[label=(\alph*)]
        \item \(x\mapsto \metric(x,A)\coloneq \inf_{z \in A} \metric(x,z)\) is \(1\)-Lipschitz continuous.
        \begin{solution}
        By the triangle inequality we have for \(x,y\in \domain\)
        \[
            \metric(x,A)
            =\inf_{z\in A} \metric(x,z)
            \le \inf_{z\in A} \bigl(\metric(x,y) + \metric(y,z)\bigr)
            = \metric(x,y) + \metric(y,A).
        \]
        with reordering this implies
        \[
            \metric(x,A) - \metric(y,A) \le \metric(x,y)
        \]
        By symmetry we also have  \(\metric(y,A) - \metric(x,A) \le \metric(x,y)\)
        and thereby the claim.
        \end{solution}

        \item \(x\mapsto x_+ = \max\set{x,0}\) is \(1\)-Lipschitz continuous.
        \begin{solution}
            Choose \(x,y\in \real\), without loss of generality
            assume \(x \ge y\). If \(y\le x\le 0\), then \(x_+ = y_+ = 0\) and thus
            \(
                \abs{x_+ - y_+} = 0 \le \abs{x-y}. 
            \)

            If \(x\ge 0\), then
            \(
                \abs{x_+ - y_+} = x - y_+ \le x-y = \abs{x-y}.
            \)
        \end{solution}

        \item If \(f\) is \(L_1\)-Lipschitz continuous and \(g\) is \(L_2\)-Lipschitz continuous, then
        \(f\circ g\) is \(L_1L_2\)-Lipschitz continuous.
        \begin{solution}
            We have for any \(x,y\in \domain\)
            \[
                \abs{f(g(x)) - f(g(y))}
                \le L_1 \abs{g(x) - g(y)}
                \le L_1L_2 \abs{x-y}.
                \qedhere
            \]
        \end{solution}


        \item the bump function \(\ind{A}^\epsilon(x) = (1 - \frac{\metric(x,A)}\epsilon)_+\)
        is \(\frac1\epsilon\)-Lipschitz continuous.
        \begin{solution}
            \(f_1(x) = \metric(x,A)\) is \(1\)-Lipschitz continuous,
            \(f_2(x) = 1 - \frac{x}{\epsilon}\) is \(\frac1\epsilon\)-Lipschitz continuous
            as a Linear function, and \(f_3(x) = x_+\) is \(1\)-Lipschitz continuous. Consequently
            \[
                f_3 \circ f_2 \circ f_1(x) = \ind{A}^\epsilon(x)
            \]
            is \(\frac1\epsilon\)-Lipschitz continuous by the previous part.
        \end{solution}
    \end{enumerate}
\end{exercise}


\begin{exercise}[How to check ``Equal in distribution'' \Coffeecup]
    \label{ex: equal in distribution}
    Let \(X\) and \(Y\) be two random variables in a metric space \(\domain\) with
    \begin{equation}
        \label{eq: equal in distribution}    
        \E{f(X)} = \E{f(Y)}
    \end{equation}
    for all bounded (measurable) functions \(f\colon \domain \to \real\). Show
    that this implies that \(X\) and \(Y\) are equal in distribution, that is
    \[
        \prob{X \in A} = \prob{Y \in A}
    \]
    for all (measurable) sets \(A\subseteq \domain\).
    
\begin{remark}
    Observe that this is very similar to Portemanteu's theorem (cf.~Definition
    \ref{def: types of convergence}) that characterizes convergence in
    distribution. And it can be shown that it is enough to check \eqref{eq:
    equal in distribution} for all bounded and Lipschitz functions \(f\colon
    \domain \to \real\).
\end{remark}
    \begin{solution}
        \(\ind{A}\) is a bounded and measurable function for any measurable set \(A\subseteq \domain\).
    \end{solution}
\end{exercise}


% \begin{exercise}[Polya's urn model]
%     In the strong law of large numbers, the random variable at the limit is a
%     degenerate random variable. A more interesting example is the Polya's urn model.
%     Consider an urn containing $R_0$ red balls and $G_0$ green balls initially. A
%     ball is drawn at random. Thus the drawn ball is red with probability
%     $R_0/(R_0+G_0)$, and is green with probability $G_0/(R_0+G_0)$. The drawn ball
%     is replaced and a ball of the same color is added to the urn. Thus if we denote
%     by $R_1$ and $G_1$ the number of red balls and green balls in the urn after
%     this, $\{R_1=R_0+1,G_1=G_0\}$ with probability $R_0/(R_0+G_0)$, and
%     $\{R_1=R_0,G_1=G_0+1\}$ with probability $G_0/(R_0+G_0)$. It can be shown that
%     as this procedure is repeated, the proportion of red balls converge almost
%     surely and the limiting random variable, say $X$, has distribution
%     $\mathrm{Beta}(R_0,G_0)$. The proportion of green balls converge to $1-X$ and
%     $1-X$ has distribution $\mathrm{Beta}(G_0,R_0)$. Thus in this situation the
%     limit is non-degenerate.
% \end{exercise}

% % \begin{lstlisting}
% % Polya <- function(N,M){
% % r = 1
% % g = 1
% % p = ggplot() +
% %   xlim(0, N) +
% %   ylim(0, 1) +
% %   xlab("n") +
% %   ylab(TeX("$\\frac{R_n}{R_n+G_n}$")) +
% %   ggtitle("Plot of Proportion of Red Balls") +
% %   theme(axis.title.y = element_text(angle = 0, vjust = 0.5)) +
% %   theme(plot.title = element_text(hjust = 0.5))
% % for(m in 1:M){
% %   g.seq = numeric(N)
% %   r.seq = numeric(N)
% %   r.seq[1] = r
% %   g.seq[1] = g
% %   for(i in 2:N){
% %     U = runif(1)
% %     if(U<r.seq[i-1]/(r.seq[i-1]+g.seq[i-1])){
% %       r.seq[i] = r.seq[i-1]+1
% %       g.seq[i] = g.seq[i-1]
% %     }
% %     else
% %     {
% %       r.seq[i] = r.seq[i-1]
% %       g.seq[i] = g.seq[i-1] + 1
% %     }
% %   }
% %   r.prop.seq = r.seq/(r.seq+g.seq)
% %   df = data.frame(x=1:N,y=r.prop.seq)
% %   p = p + geom_line(data=df, aes(x=x,y=y), color=m)
% %   }
% %   return(p)
% % }
% % \end{lstlisting}

% \begin{figure}[h]
%     \centering
%     % \includegraphics[width=\linewidth]{Figures/Rplot02.png}
%     \caption{Plot of $R_n/(R_n+G_n)$ for $n=1,2,\ldots,1000$. Here $R_0=1,G_0=1$. $10$ simulations are shown. In all cases the proportions approach a limit, and the limiting random variable
%     is non-degenerate. The limit is clearly different for the different simulations.}
%     \label{fig:enter-label}
% \end{figure}

